\documentclass[border=4pt]{standalone}
\usepackage{amsmath,amssymb,mathtools,xcolor,varwidth}
\begin{document}
\begin{varwidth}{28em}
\large\bfseries
A distance travelled is velocity accumulated over time, and that accumulation is precisely      the AREA swept beneath the velocity curve. So the areas $\triangle ABD$ and      $\triangle ACE$ — the spaces described in the two intervals are what we must compare. But      this is the identical figure of Lemma IX, and by Lemma IX these areas are ultimately in the      duplicate ratio of the base sides, namely $AD^2 : AE^2$. Since $AD$ and $AE$ are      the times, the spaces stand as the squares of the times — the celebrated law $s \propto t^2$      that governs everything falling from rest. Cor.\ 4 then closes the ring: rearranging      $s \propto F\,t^2$ we read off $F \propto \dfrac{s}{t^2}$, so the forces are as the spaces      directly and the squares of the times inversely. \textit{Q.E.D.}
\end{varwidth}
\end{document}
